\section{Limitations}
\label{sec:limitations}

The independent FY2023--24 estimate covers four IPPU categories: cement calcination
(2.A.1, \SI{21.815648}{\mega\tonne}), ammonia production (2.B.1, \SI{3.2433}{\mega\tonne}),
iron and steel (2.C.1, \SI{0.3931}{\mega\tonne}) and lubricant use (2.D.1,
\SI{0.0833}{\mega\tonne}), for a module total of \SI{25.5354}{\mega\tonne} of \COtwo{}.
Methods and data requirements are defined for the categories that remain unestimated --- lime
(2.A.2), glass (2.A.3), nitric acid (2.B.2), HFCs (2.F) and SF$_6$ (2.G) --- and those stay
NE until the required activity and process data are available. Cement is therefore
\SI{85.4}{\percent} of the module rather than all of it: the
\SI{21.815648}{\mega\tonne} figure represents cement-process \COtwo{} rather than total
IPPU or total cement emissions, and the \SI{17.389}{\mega\tonne} Tier~1 inference at the 0.75
clinker fraction is a named diagnostic beside it, not the reference.

National clinker production is observed at the company grain; stock change is not. All sixteen
producer records together exceed the low-clinker inference by a quarter and the high case by
1.4\%, but the company-level purchase, sale, grinding, export and stock fields remain
missing. The company compilation therefore constrains the national clinker balance but does
not yet constitute a complete Tier~2 activity dataset.

The conditional Monte Carlo covers activity, trade and factor ranges within each clinker-
fraction model. It does not assign a probability to the 0.75 or 0.95 configuration. Stock
uncertainty is not sampled because no defensible public bound was identified.

The matched calendar-2022 comparison still depends on an assumed clinker fraction and zero
stock change. It demonstrates that year mismatch does not explain the difference from the CRT,
but it cannot reconstruct the official institutional derivation without its internal clinker
and coverage evidence.

The chemical plant table is partial. Nameplate capacity, one calendar-2024 ammonia observation and CDM
project existence cannot substitute for national plant production, recovered \COtwo{} or
annual abatement operation. Ozone-profile HFC consumption provides a national supply statistic
but cannot be compared directly with annual emissions.

Company-level clinker production does not establish plant-level output. Facility-assisted geography uses only the source locations compatible with the primary package's licence; ODbL-derived alternatives remain separate. Unsupported or excluded locations remain unallocated rather than inheriting another source's coordinates.

The common annual core covers ten fiscal years, July 2015--June 2025. The reference-year company analysis does not imply observed annual plant geography. Soda ash production is observed; its Solvay-route GHG treatment places the coke combustion term under Energy, where missing coke activity remains explicit. The IPPU notation must not be read as evidence that soda ash production does not occur. The admitted soda-ash air process class has its own factor and boundary. Lime, glass, nitric acid and fluorinated-gas emissions remain unresolved where required process or equipment data are missing.
